\documentclass[11pt,a4paper]{article}
\usepackage{fontspec}
\usepackage[margin=1in]{geometry}
\usepackage[unicode,hidelinks]{hyperref}
\setlength{\parindent}{0pt}
\setlength{\parskip}{.6em}
\emergencystretch=3em
\newcommand{\volumeentry}[3]{\par\begingroup\parskip=.7em
  \noindent\href{TDOME_Complete_Collection.pdf\#page=#3}{\textbf{#1}\quad #2\nobreak\dotfill\nobreak\mbox{#3}}\par\endgroup}
\hypersetup{pdftitle={Thermodynamic Dynamics of Observer-Memory Entanglement},pdfauthor={Sidong Liu}}
% COVER SUBTITLE: Memory, self-frames and continued calibration
\begin{document}
\begin{titlepage}
\centering
\vspace*{1.6cm}
{\Huge\bfseries Thermodynamic Dynamics of Observer-Memory Entanglement\par}
\vspace{1cm}
{\Large Memory, self-frames and continued calibration\par}
\vfill
{\Large Complete Collected Volume\par}
\vspace{1.6cm}
{\Large Sidong Liu, PhD\par}
\vspace{.3cm}
{iBioStratix Ltd\par}
\href{mailto:sidongliu@hotmail.com}{\texttt{sidongliu@hotmail.com}}\par
\vspace{1cm}
{\large October 2026\par}
\end{titlepage}
\setcounter{page}{2}
\section*{Contents}
The numbers below refer to pages in this PDF. Each paper retains its own
printed pagination, complete proofs, appendices and references.
\volumeentry{}{Introduction and reading map}{3}
\volumeentry{TD-I}{Memory as a Usable Resource: Correlations, Records and Complete Endpoint Accounting}{4}
\volumeentry{TD-II}{Reference Frames Selected by History: The Capability and Cost of a Bounded Agent's Self-Frame}{27}
\volumeentry{TD-III}{Keeping a Self-Model Calibrated: Drift Detection, the Price of Keeping Up, and Reliability in a Closed Loop}{47}
\clearpage
\section*{Introduction}
% COLLECTION INTRODUCTION START

An agent need not acquire a complete image of itself to keep acting. It needs records that improve the decisions available to its continuing body. The three papers develop memory as a resource, a history-selected frame, and continued calibration when that frame drifts or the body loses access. The central claim concerns functional necessity in specified tasks, not an unconditional theorem that every memory generates experience.

I supplies the correlation and thermodynamic account. II constructs the usable frame. III determines when stored history changes detection, maintenance and action decisions. Reading them in this order exposes what each transition requires, including paid references and information that must remain inaccessible in a claimed decision witness.



\section*{Reading map}
\begin{description}
\item[TD-I] \textit{Memory as a Usable Resource: Correlations, Records and Complete Endpoint Accounting}. The work value of records depends on their physical readout and retained side information. The endpoint account keeps input purity and the outputs of correction visible.
\item[TD-II] \textit{Reference Frames Selected by History: The Capability and Cost of a Bounded Agent's Self-Frame}. An agent's usable frame is selected by its actual history. Oriented records and legal readouts determine what it can calibrate and control.
\item[TD-III] \textit{Keeping a Self-Model Calibrated: Drift Detection, the Price of Keeping Up, and Reliability in a Closed Loop}. Drift makes continued self-calibration a task. Retained histories determine reference requests, paid maintenance and action authority on one continuing body.
\end{description}

The papers retain their complete proofs, appendices and references, and their own pagination. The contents list PDF page numbers; PDF bookmarks also provide section-level navigation.
% COLLECTION INTRODUCTION END
\end{document}
